% Part: sets-functions-relations
% Chapter: size-of-sets
% Section: pairing
% Hindi translation (hi-Deva-IN) of Open Logic commit 9620cc73f9c8e0ad003c514a5d3748f29611c4c0.

\documentclass[../../../include/open-logic-section]{subfiles}

\begin{document}

\olfileid[hi]{sfr}{siz}{pai}
\olsection{युग्मन फलन और कूट}

\begin{explain}
कैंटर की आड़ी-तिरछी विधि से $\Nat^n$ की गणनीयता दृष्टिगत रूप से स्पष्ट हो
जाती है। किंतु आइए, $\Nat^2$ को दर्शाने वाली अपनी सरणी पर ध्यान केंद्रित करें।
सरणी में आड़ी-तिरछी रेखा के साथ चलते हुए और स्थानों को गिनते हुए हम जाँच सकते
हैं कि $\tuple{1,2}$ संख्या~$7$ से संबद्ध है। फिर भी, अच्छा होता यदि हम इसे
और सीधे ढंग से परिकलित कर पाते। अर्थात, अच्छा होता यदि हमारे पास आड़ी-तिरछी
गणना का \emph{प्रतिलोम} $g\colon \Nat^2 \to \Nat$ होता, जहाँ
\[
g(\tuple{0,0}) = 0, \;
g(\tuple{0,1}) = 1, \;
g(\tuple{1,0}) = 2, \; \dots,  
g(\tuple{1,2}) = 7, \; \dots
\]
इससे हम ठीक-ठीक परिकलित कर सकते कि हमारी गणना में $\tuple{n, m}$ कहाँ आएगा।

वस्तुतः, दो बातों पर ध्यान देकर हम $g$ को सीधे परिभाषित कर सकते हैं। पहली:
यदि $n$वीं पंक्ति और $m$वें स्तंभ में मान~$v$ है, तो $(n+1)$वीं पंक्ति और
$(m-1)$वें स्तंभ में मान $v + 1$ है। दूसरी: हमारी गणना की पहली पंक्ति में
$0$, $1$, $3$, $6$ इत्यादि से आरंभ होने वाली त्रिकोणीय संख्याएँ हैं। $k$वीं
त्रिकोणीय संख्या $< k$ प्राकृत संख्याओं का योग है, जिसे $k(k+1)/2$ के रूप में
परिकलित किया जा सकता है। इन दोनों बातों को एक साथ रखकर इस फलन पर विचार करें:
\[
  g(n,m) = \frac{(n+m+1)(n+m)}{2} + n
\]
हम प्रायः केवल $g(n, m)$ लिखते हैं, न कि $g(\tuple{n, m})$, क्योंकि इसे
पढ़ना अधिक सुगम है। यह हमें पहले $(n+m)^\text{वीं}$ त्रिकोणीय संख्या निर्धारित करने
और फिर उसमें $n$ जोड़ने को कहता है। इससे सरणी ठीक उसी प्रकार भरती है जैसा हम
चाहते हैं। अतः, विशेष रूप से, युग्म $\tuple{1, 2}$ को
$\frac{4 \times 3}{2} +
1 = 7$ पर भेजा जाता है।

यह फलन $g$ युग्मों के किसी समुच्चय की गणना का \emph{प्रतिलोम} है। ऐसे फलन
\emph{युग्मन फलन} कहलाते हैं।
\end{explain}

\begin{defn}[युग्मन फलन]
  कोई फलन $f\colon A \times B \to \Nat$ एक अंकगणितीय \emph{युग्मन फलन}
  है, यदि $f$ अंतःक्षेपी है। हम यह भी कहते हैं कि $f$, $A \times B$ को
  \emph{कूटित करता है,} और $f(x,y)$, $\tuple{x,y}$ का \emph{कूट} है।
\end{defn}

\begin{explain}
हम युग्मन फलनों का उपयोग, उदाहरणार्थ, प्राकृत संख्याओं के युग्मों को कूटित
करने में कर सकते हैं; दूसरे शब्दों में, हम अवयवों के प्रत्येक \emph{युग्म}
को किसी \emph{एक} संख्या द्वारा निरूपित कर सकते हैं। युग्मन फलन के प्रतिलोम
का उपयोग करके हम उस संख्या को \emph{विकूटित} कर सकते हैं, अर्थात पता लगा
सकते हैं कि वह किस युग्म का निरूपण करती है।
\end{explain}

\begin{prob}
सभी गैर-ऋणात्मक परिमेय संख्याओं के समुच्चय की एक गणना दीजिए।
\end{prob}

\begin{prob}
दिखाइए कि $\Rat$ !!{enumerable} है। याद कीजिए कि किसी भी परिमेय संख्या को
$z/m$ भिन्न के रूप में लिखा जा सकता है, जहाँ $z \in \Int$, $m \in \Nat^+$।
\end{prob}

\begin{prob}
$\Bin^*$ की एक गणना परिभाषित कीजिए।
\end{prob}

\begin{prob}
अपने परिचयात्मक तर्कशास्त्र पाठ्यक्रम से याद कीजिए कि प्रत्येक संभव सत्यमान
सारणी एक सत्यमान फलन व्यक्त करती है। दूसरे शब्दों में, सत्यमान फलन वे सभी
$\Bin^k \to \Bin$ फलन हैं, जहाँ ऐसा कोई~$k$ है। सिद्ध कीजिए कि सभी
सत्यमान फलनों का समुच्चय गणनीय है।
\end{prob}

\begin{prob}
दिखाइए कि किसी भी अनंत !!{enumerable} समुच्चय के सभी परिमित उपसमुच्चयों
का समुच्चय !!{enumerable} है।
\end{prob}

\begin{prob}
$\Nat$ का कोई उपसमुच्चय \emph{सहपरिमित} कहलाता है यदि और केवल यदि वह
$\Nat$ में किसी परिमित समुच्चय का पूरक हो; अर्थात, $A \subseteq \Nat$
सहपरिमित है यदि और केवल यदि $\Nat\setminus A$ परिमित है। मान लीजिए $I$ वह
समुच्चय है जिसके !!{element} ठीक $\Nat$ के परिमित और सहपरिमित उपसमुच्चय हैं।
दिखाइए कि $I$ !!{enumerable} है।
\end{prob}

\begin{prob}
दिखाइए कि !!{enumerable} समुच्चयों का !!{enumerable} संघ !!{enumerable}
होता है। अर्थात, जब भी $A_1$, $A_2$, \dots{} समुच्चय हों और प्रत्येक $A_i$
!!{enumerable} हो, तब उन सभी का संघ $\bigcup_{i=1}^\infty
A_i$ भी
!!{enumerable} होता है। [ध्यान दें: यह कठिन है!]
\end{prob}

\begin{prob}
मान लीजिए $f \colon A \times B \to \Nat$ कोई भी युग्मन फलन है।
दिखाइए कि $f$ का प्रतिलोम $A \times B$ की एक गणना है।
\end{prob}

\begin{prob}
$\Nat^3$ को कूटित करने वाला कोई फलन निर्दिष्ट कीजिए।
\end{prob}

\end{document}
